\documentclass{article}
\usepackage{amsmath, amssymb, amsthm}

\title{IMC 2005, Blagoevgrad \\ First Day}
\date{}

\begin{document}
\maketitle

\section*{Problem 1}
Let $A$ be the $n \times n$ matrix with $(i, j)$-entry $i + j$. Find $\operatorname{rank}(A)$.

\subsection*{Solution}
$A_{ij} = i + j$, so $A = (i)_{ij} + (j)_{ij}$ is a sum of two rank-$1$ matrices, hence $\operatorname{rank}(A) \le 2$. For $n \ge 2$, the top-left $2 \times 2$ minor has determinant $2 \cdot 4 - 3 \cdot 3 = -1 \ne 0$, so $\operatorname{rank}(A) = 2$. For $n = 1$, $\operatorname{rank}(A) = 1$.

\section*{Problem 2}
For $n \ge 3$, let
\[
S_n = \{0, 1, 2\}^n, \quad A_n = \{x \in S_n : |\{x_i, x_{i+1}, x_{i+2}\}| \ne 1 \text{ for all } i \le n-2\},
\]
\[
B_n = \{x \in S_n : x_i = x_{i+1} \Rightarrow x_i \ne 0\}.
\]
Prove $|A_{n+1}| = 3 |B_n|$.

\subsection*{Solution}
Cyclic shift $x \mapsto x + (1, 1, \ldots, 1) \pmod 3$ acts on $A_n$, so $1/3$ of $A_{n+1}$ starts with $0$. Define a bijection:
\[
(0, x_1, \ldots, x_n) \in A_{n+1} \longleftrightarrow (y_1, \ldots, y_n) \in B_n, \quad y_k = x_k - x_{k-1} \pmod 3 \text{ (with } x_0 = 0\text{)}.
\]
If $y_k = y_{k+1} = 0$ (in $\mathbb{Z}_3$), then $x_{k-1} = x_k = x_{k+1}$ would violate the $A$ condition. Inverse: $x_k = y_1 + \cdots + y_k$.

\section*{Problem 3}
$f : \mathbb{R} \to [0, \infty)$ is $C^1$. Prove
\[
\Bigl| \int_0^1 f^3 - f^2(0) \int_0^1 f \Bigr| \le \max_{[0,1]} |f'| \cdot \Bigl(\int_0^1 f\Bigr)^2.
\]

\subsection*{Solution}
Let $M = \max|f'|$. From $-M \le f' \le M$: $-M f \le f f' \le M f$. Integrating from $0$ to $x$: $-M \int_0^x f \le \tfrac{1}{2}(f^2(x) - f^2(0)) \le M \int_0^x f$. Multiply by $f(x) \ge 0$:
\[
-M f(x) \int_0^x f \le \tfrac{1}{2} f^3(x) - \tfrac{1}{2} f^2(0) f(x) \le M f(x) \int_0^x f.
\]
Integrate over $[0, 1]$; using $\int_0^1 f(x) \int_0^x f(t) dt dx = \tfrac{1}{2}(\int_0^1 f)^2$:
\[
\Bigl|\int_0^1 f^3 - f^2(0) \int_0^1 f\Bigr| \le M (\int_0^1 f)^2.
\]

\section*{Problem 4}
Find all polynomials $P(x) = \sum_{i=0}^n a_i x^i$ ($a_n \ne 0$) with (i) $(a_0, \ldots, a_n)$ a permutation of $(0, 1, \ldots, n)$; (ii) all roots rational.

\subsection*{Solution}
All roots are nonpositive (since $P(x) > 0$ for $x > 0$). Write roots as $-\alpha_i, \alpha_i \ge 0$.

Must have $a_0 = 0$: else if some $a_k = 0$ with $k \ge 1$, Vieta gives a sum of products of $\alpha_i$'s equals 0, impossible for positive $\alpha_i$. So one root is $0$; factor out $x$.

For degree $\ge 3$, from Vieta and AM-HM on the remaining $n - 1$ roots: $\tfrac{a_{n-1}}{(n-1) a_n} \ge \tfrac{(n-1) a_1}{a_2}$, giving $\tfrac{a_2 a_{n-1}}{a_1 a_n} \ge (n-1)^2$. Since LHS $\le n^2/2$, we need $n \le 3$.

Checking: $P(1) = n(n+1)/2 = \prod(b_k + c_k) \ge 2^{n-1}$ (when factored in $\mathbb{Z}[x]$), giving $n \le 4$; case $n = 4$ rules out by factorization.

Complete list: $x$, $x^2 + 2x$, $2x^2 + x$, $x^3 + 3x^2 + 2x$, $2x^3 + 3x^2 + x$.

\section*{Problem 5}
$f \in C^2((0, \infty))$ with $|f''(x) + 2x f'(x) + (x^2 + 1) f(x)| \le 1$. Prove $\lim_{x \to \infty} f(x) = 0$.

\subsection*{Solution}
Let $g = f' + xf$; then $f'' + 2xf' + (x^2 + 1)f = g' + xg$.

\textbf{Lemma:} If $h$ is $C^1$ and $|h' + xh| \le M$, then $h(x) \to 0$.

Let $p(x) = h(x) e^{x^2/2}$. Then $|p'(x)| = |h' + xh| e^{x^2/2} \le M e^{x^2/2}$. So
\[
|h(x)| = \frac{|p(x)|}{e^{x^2/2}} \le \frac{|p(0)| + M \int_0^x e^{t^2/2} dt}{e^{x^2/2}} \to 0
\]
by L'Hôpital ($\tfrac{d}{dx} \int e^{t^2/2}dt / \tfrac{d}{dx} e^{x^2/2} = 1/x$).

Apply the lemma to $g$, then to $f$.

\section*{Problem 6}
$G$ group; $G(m)$ = subgroup generated by $m$-th powers. If $G(m)$ and $G(n)$ are commutative, prove $G(\gcd(m, n))$ is commutative.

\subsection*{Solution}
Let $d = \gcd(m, n)$. Then $\langle G(m), G(n) \rangle = G(d)$ (from $d = um + vn$: $g^d = (g^m)^u (g^n)^v$ under appropriate grouping; more carefully, $G(d)$ contains $G(m)$ and $G(n)$, and $g^d \in \langle G(m), G(n)\rangle$).

Suffices to show two generators commute: $a^m$ and $b^n$. Let $z = a^{-m} b^{-n} a^m b^n$. Then
\[
z = (a^{-m} b a^m)^{-n} b^n = a^{-m} (b^{-n} a b^n)^m,
\]
so $z \in G(m) \cap G(n)$. Hence $z$ is central in both $G(m)$ and $G(n)$, so central in $G(d)$.

From $a^m b^n = b^n a^m z$, induct: $a^{ml} b^{nl} = b^{nl} a^{ml} z^{l^2}$. Set $l = m/d$: since $(b^n)^{m/d} = b^{mn/d} \in G(m)$, and $G(m)$ commutative, conjugation of $a^m$ by $b^{mn/d}$ is trivial, forcing $z^{(m/d)^2} = e$. Similarly $z^{(n/d)^2} = e$. Since $\gcd(m/d, n/d) = 1$, $z = e$.

\end{document}
